\documentclass[10pt,a4paper]{article}
\usepackage[a4paper,margin=17mm,top=16mm,bottom=18mm]{geometry}
\usepackage{fontspec,xcolor,graphicx,tabularx,enumitem,hyperref,xurl,fancyhdr,titlesec,needspace,microtype}
\usepackage[ngerman]{babel}
\setlength{\emergencystretch}{2em}
\setmainfont{Lato}\setsansfont{Lato}\renewcommand{\familydefault}{\sfdefault}
\definecolor{Ink}{HTML}{19373F}\definecolor{Teal}{HTML}{176678}\definecolor{Muted}{HTML}{4C636B}\definecolor{Rust}{HTML}{A85731}\definecolor{Rule}{HTML}{D3DFDF}
\hypersetup{colorlinks=true,urlcolor=Teal,linkcolor=Teal,pdfauthor={Dr. Chilperic Armel Foko Kuate},pdftitle={Curriculum Vitae — Dr. Chilperic Armel Foko Kuate}}
\setlength{\parindent}{0pt}\setlength{\parskip}{3pt}\linespread{1.06}\color{Ink}
\setlist[itemize]{leftmargin=1.25em,itemsep=2pt,topsep=3pt,parsep=0pt,label={\color{Rust}\small\textbullet}}
\pagestyle{fancy}\fancyhf{}\renewcommand{\headrulewidth}{0pt}\renewcommand{\footrulewidth}{.4pt}\fancyfoot[L]{\scriptsize\color{Muted}Dr. Chilperic Armel Foko Kuate}\fancyfoot[C]{\scriptsize\color{Muted}Aktualisiert: 4. Oktober 2026}\fancyfoot[R]{\scriptsize\color{Teal}\thepage}
\titleformat{\section}{\large\bfseries\color{Teal}}{}{0pt}{}[\vspace{2pt}{\color{Rule}\titlerule}]\titlespacing*{\section}{0pt}{10pt}{5pt}
\newcommand{\role}[3]{\Needspace{5\baselineskip}{\bfseries #1}\hfill{\small\bfseries\color{Teal}#2}\par{\small\color{Muted}#3}\par}
\newcommand{\degree}[3]{\Needspace{3\baselineskip}{\bfseries #1}\hfill{\small\bfseries\color{Teal}#2}\par{\small\color{Muted}#3}\par\vspace{2pt}}
\newcommand{\project}[2]{\Needspace{4\baselineskip}{\bfseries #1}\par #2\par\vspace{3pt}}
\newcommand{\continued}[1]{{\small\bfseries\color{Teal}DR. CHILPERIC ARMEL FOKO KUATE}\hfill{\small\color{Muted}#1}\par\vspace{5pt}}
\newcommand{\headline}[2]{{\small\bfseries\color{Rust}#1}\par\vspace{5pt}{\fontsize{27}{29}\selectfont\bfseries Dr. Chilperic Armel\\Foko Kuate}\par\vspace{7pt}{\large\color{Teal}#2}\par\vspace{9pt}{\small Düsseldorf, Germany\quad\textbar\quad\href{mailto:chilpericarmel@gmail.com}{chilpericarmel@gmail.com}}\par{\small\href{https://chilperic.github.io}{Research website}\quad\textbar\quad\href{https://github.com/chilperic}{GitHub}\quad\textbar\quad\href{https://orcid.org/0000-0002-0140-7588}{ORCID: 0000-0002-0140-7588}}\par\vspace{7pt}{\color{Rust}\rule{23mm}{1.5pt}}{\color{Rule}\rule{\dimexpr\linewidth-23mm}{.5pt}}\par}
\begin{document}
\headline{WISSENSCHAFTLICHE SOFTWARE \& COMPUTATIONAL SCIENCE}{Mechanistische Modellierung · Numerik · Datenanalyse}
\section{Profil}
Computational Scientist mit Promotion in quantitativer und theoretischer Biologie. Forschungserfahrung in mechanistischer Modellierung, nichtlinearer Dynamik und wissenschaftlicher Software. Ich entwickle mathematische Modelle und reproduzierbare Arbeitsabläufe, die Annahmen, Simulationen, Daten und Interpretation verbinden.
\section{Fachliche Schwerpunkte}
\begin{tabularx}{\linewidth}{@{}p{35mm}X@{}}
\textbf{Modellierung} & ODE-/DAE-Systeme, Enzymkinetik, stochastische Prozesse, Zellpopulationen, Gleichgewichte und Stabilität.\\[5pt]
\textbf{Analyse} & Parameterschätzung, Sensitivitäts- und Unsicherheitsanalyse, beschränkte und ableitungsfreie Optimierung, statistische Datenanalyse.\\[5pt]
\textbf{Software} & Python, NumPy/SciPy, SALib, CasADi/IPOPT, CMA-ES; Git, Linux, Jupyter, Tests und Dokumentation. Erfahrung mit Julia, R, MATLAB und wissenschaftlichen JavaScript-Oberflächen.
\end{tabularx}
\section{Berufliche Erfahrung}
\role{Wissenschaftlicher Mitarbeiter · Computational Cell Biology}{seit Okt 2026}{Heinrich-Heine-Universität Düsseldorf · CCB}
Wiederaufnahme der Forschungstätigkeit am 1. Oktober 2026; rechnergestützte Biologie und mechanistische Modellierung.
\role{Eigenständige Forschung \& FokoLab-Entwicklung}{Nov 2025–Sep 2026}{Düsseldorf, Deutschland}
Entwicklung browserbasierter wissenschaftlicher Werkzeuge, Pflege von Modellierungsprojekten, wissenschaftliches Schreiben und Weiterbildung in KI/ML.
\role{Postdoktorand}{Nov 2022–Okt 2025}{CEPLAS / CCB, Heinrich-Heine-Universität Düsseldorf}
\begin{itemize}
\item Verknüpfung biochemischer, thermischer und hydraulischer Prozesse in Modellen zur Anpassung von Pflanzen; Analyse von Umweltreaktionen und physiologischen Zielkonflikten.
\item Entwicklung von Simulations-, Optimierungs- und Sensitivitätsanalysen sowie Betreuung studentischer Forschungsprojekte.
\end{itemize}
\role{Doktorand · PoLiMeR ITN}{Aug 2019–Aug 2022}{Quantitative und Theoretische Biologie, HHU Düsseldorf}
Kinetische Modelle der Fettsäuresynthese und des hepatischen Stoffwechsels; Gleichgewichts- und Stabilitätsanalysen; Beiträge zu begutachteten Veröffentlichungen.
\newpage\continued{Projekte \& Qualifikationen}
\role{Research Fellow}{Nov 2017–Mär 2019}{AIMS · Kamerun und Ruanda}
Deterministische und stochastische Modelle der T-Zell-Proliferation; Lehre und wissenschaftliches Rechnen.
\section{Ausgewählte Projekte}
\project{FokoLab · wissenschaftliche Anwendung}{Verbundene Labore für Simulation, Optimierung, Sensitivität, Parameterschätzung, Statistik und mechanistische Visualisierung. Eigene Modelle, gespeicherte Experimente und reproduzierbare Exporte.\\\href{https://chilperic.github.io/chilperic_ode_solver/site/index.html}{Webanwendung} · \href{https://github.com/chilperic/chilperic_ode_solver}{Quellcode}}
\project{Pflanzen–Umwelt-Modellierung}{Biochemisch-thermisch-hydraulischer Ansatz zur photosynthetischen Anpassung; reproduzierbare Szenarien und Zusammenarbeit mit Studierenden und experimentell Forschenden.\\\href{https://github.com/chilperic/photosynthesis-climate-adaptation-model}{Forschungssoftware}}
\project{Stoffwechsel- und Zellpopulationsmodelle}{Fettsäurekinetik, nichtlineare Stoffwechseldynamik, Verzweigungsprozesse und Modelle von Zellteilungsdaten.\\\href{https://github.com/chilperic}{Weitere Repositorien}}
\section{Ausbildung}
\degree{PhD · Quantitative und Theoretische Biologie}{2019–2023}{Heinrich-Heine-Universität Düsseldorf. Abschluss: 7. März 2023.}
\degree{MSc · Mathematik und grundlegende Anwendungen}{2017–2019}{Schwerpunkt Modellierung · Universität Yaoundé I, Kamerun.}
\degree{MSc · Mathematical Sciences}{2015–2016}{AIMS Ghana / University of Cape Coast, Ghana.}
\degree{BSc · Mathematik}{2009–2012}{Universität Douala, Kamerun.}
\degree{Baccalauréat}{2009}{Collège La Confiance, Bafoussam, Kamerun.}
\section{Ausgewählte Veröffentlichungen}
{\small Foko Kuate, Ebenhöh, Bakker \& Raguin (2023). Enzymkinetische Daten zur tierischen Fettsäuresynthese. \textit{Bioscience Reports}. \href{https://doi.org/10.1042/BSR20222496}{\nolinkurl{10.1042/BSR20222496}}.\\[4pt]
Wilken et al., einschließlich Foko Kuate (2022). Differenzierbare, beschränkungsbasierte Stoffwechselmodelle. \textit{Metabolic Engineering} 74, 72–82. \href{https://doi.org/10.1016/j.ymben.2022.09.002}{\nolinkurl{10.1016/j.ymben.2022.09.002}}.}
\section{Lehre, Sprachen \& Weiterbildung}
Betreuung von zwei MSc- und zwei BSc-Projekten; über 200 Lehr-/Übungsstunden. Französisch: Muttersprache; Englisch: gute berufliche Kenntnisse; Deutsch: B1/B1+, B2-Weiterbildung laufend. Weiterbildung in KI-Engineering und Machine Learning.

\end{document}
